\chapter{Data Integrity Audits}
\label{ch:appendix_a}

This appendix records the systematic audit of data-layer assumptions conducted before the analysis was finalised.
It documents what was checked, what was found, and how each issue was addressed.
The main text reports only the corrected system and results, while this appendix retains the revision trail needed to reproduce the analysis and reconcile differences with earlier project-stage documents, including the interim presentation.

\begin{longtblr}[
		caption = {Audit findings and corrective actions.},
		label = {tab:audit_findings},
	]{
		width = \linewidth,
		colspec = {Q[l,m,2.6cm] X[l] X[l]},
		colsep = 4pt,
		row{1} = {font=\bfseries},
		rowhead = 1,
		rowsep = 2pt,
	}
	\toprule
	Audit item                    & Finding                                                                                                                                                                                                     & Corrective action                                                                                                                                                                                                                                                \\
	\midrule
	Feature provenance            & Compactness and floor area were certificate fields used to calculate the target. Their inclusion as predictors constituted target leakage, and the fields would not be available for buildings without EPCs. & Removed both fields, replaced them with reconstructed 3DBAG-derived values, and updated all affected results.                                                                                                                                                     \\
	TABULA lookup                 & The manually transcribed table contained a systematic thermal-resistance bias relative to the source workbook.                                                                                               & Replaced manual transcription with automated generation and value-level provenance. Classification results were unaffected because the lookup values are deterministic functions of type and period; physical quantities were recalculated using corrected values. \\
	Target definition             & The registered energy class is derived from \texttt{Primaire\allowbreak Fossiele\allowbreak Energie\allowbreak EMG\allowbreak Forfaitair} where that field is available, rather than from the unadjusted primary fossil energy field. & Aligned the regression target with the registered energy class definition and recalculated the affected regression results.                                                                                                       \\
	Count reconciliation          & Building counts reported at different stages referred to different data populations.                                                                                                                        & Standardised the convention of reporting 10,104 buildings in the filtered \ac{SVI} dataset and 10,086 buildings in the experimental dataset.                                                                                                                      \\
	Certificate chronology        & The database contains only current certificates, and only 0.1\% of dwellings nationwide have a certificate time series.                                                                                      & The available data did not support longitudinal renovation labels; renovation detection was retained as future work in \cref{sec:future_work}.                                                                                                                    \\
	Nature of the geometry source & 3DBAG contains automatically reconstructed outputs rather than field measurements. Floor count and the party-wall and external-wall areas are derived fields introduced after \textcite{peters2022}.          & Standardised the manuscript terminology as \emph{reconstructed geometry}, documented the derivation of each field in \cref{tab:3dbag_fields}, and identified the 3DBAG version as required metadata.                                                               \\
	\bottomrule
\end{longtblr}

The feature-provenance finding was converted into the general eligibility rule stated in \cref{sec:downstream_model}: the provenance of every candidate predictor must be verified, and any field used to calculate the target must be excluded.

\begin{table}[htbp]
	\centering
	\caption{Definitions of derived 3DBAG fields referenced in \cref{sec:data_sources}.}
	\label{tab:3dbag_fields}
	\footnotesize
	\begin{tblr}{
			width = \linewidth,
			colspec = {Q[l,m,3.4cm] X[l] X[l]},
			colsep = 4pt,
			row{1} = {font=\bfseries},
			rowsep = 2pt,
		}
		\toprule
		Field used in this study        & Definition                                                                    & Note                                                                     \\
		\midrule
		\texttt{building\_height}       & \texttt{b3\_h\_max} $-$ \texttt{b3\_h\_maaiveld}                              & Ground elevation is the fifth percentile of ground-point elevations.     \\
		\texttt{num\_floors\_estimated} & \texttt{b3\_bouwlagen}; if missing, \texttt{round(building\_height / 3.0)}     & Assumes a floor-to-floor height of 3.0\,m.                               \\
		\texttt{floor\_area\_estimated} & \texttt{b3\_opp\_grond} $\times$ \texttt{num\_floors\_estimated}               & Replaces the certificate field used incorrectly in the earlier analysis. \\
		\texttt{shared\_ratio}          & \texttt{b3\_opp\_scheidingsmuur / (b3\_opp\_buitenmuur + b3\_opp\_scheidingsmuur)} & Classification statistic derived from reconstructed surfaces.        \\
		\texttt{shape\_factor}          & \texttt{envelope\_area / volume}                                              & A comparison variant that includes shared walls is also retained.        \\
		\bottomrule
	\end{tblr}
	\par\vspace{0.4em}
	\begin{minipage}{\linewidth}
		% TODO (author action): insert the identifier of the 3DBAG version used.
		\footnotesize\textit{To be completed:} identifier of the 3DBAG version used.
	\end{minipage}
\end{table}

\begin{table}[htbp]
	\centering
	\caption[Sensitivity of the $h_{tr}$ MAE to the assumed WWR.]{Sensitivity of the $h_{tr}$ MAE ($\mathrm{W/(K\,m^2)}$) to WWR (\cref{sec:weighted_archetype_metric}; corresponding to \cref{tab:htr_metrics}).}
	\label{tab:wwr_sensitivity}
	\small
	\begin{tblr}{
			width = \linewidth,
			colspec = {X[l] Q[r,m,2.4cm] Q[r,m,2.4cm] Q[r,m,2.4cm]},
			colsep = 4pt,
			row{1} = {font=\bfseries},
			rowsep = 2pt,
		}
		\toprule
		Vision configuration   & WWR 0.15 & WWR 0.25 & WWR 0.35 \\
		\midrule
		DINOv2 (frozen)        & 0.150    & 0.150    & 0.150    \\
		ResNet-50 (fine-tuned) & 0.157    & 0.157    & 0.157    \\
		InternVL3 (zero-shot)  & 0.495    & 0.474    & 0.453    \\
		\bottomrule
	\end{tblr}
\end{table}

The ranking and separation of the three configurations remained unchanged across the tested WWR values.
The MAEs of the two trained configurations were unchanged at the reported precision because their cell errors were concentrated among the adjacent early periods NL.01--NL.03, which share the same window U-value of 5.2\,$\mathrm{W/(m^2K)}$ (\cref{tab:tabula_uvalues}).
The smaller variation for InternVL3-2B is consistent with assignments that cross the change in window U-value from 5.2 to 1.8\,$\mathrm{W/(m^2K)}$.
Source: \texttt{reports/tables/stage3/\allowbreak T7\_htr\_instrument.md}, generated by \texttt{src/stage3/htr\_instrument.py}.

\begin{table}[htbp]
	\centering
	\caption[Per-class classification report for the binary task.]{Per-class classification report for the binary task (\cref{tab:binary_routes}; holdout set, $n = 2{,}018$; support for classes A to C = 1,416; support for classes D to G = 602).}
	\label{tab:binary_perclass}
	\footnotesize
	\begin{tblr}{
			width = \linewidth,
			colspec = {Q[l,m,2.4cm] Q[r,m,1.85cm] Q[r,m,1.85cm] Q[r,m,1.85cm] Q[r,m,1.85cm] Q[r,m,1.85cm] Q[r,m,1.85cm]},
			colsep = 2pt,
			row{1} = {font=\bfseries},
			rowsep = 2pt,
		}
		\toprule
		Route                 & A-C P        & A-C R        & A-C F1       & D-G P        & D-G R        & D-G F1       \\
		\midrule
		M0 (random, expected) & 0.702          & 0.500          & 0.584          & 0.298          & 0.500          & 0.373          \\
		M1                    & 0.715          & \textbf{0.960} & \textbf{0.819} & \textbf{0.509} & 0.098          & 0.165          \\
		M2-DINOv2            & \textbf{0.764} & 0.731          & 0.748          & 0.426          & 0.469          & \textbf{0.447} \\
		M2-ResNet-50         & 0.741          & 0.741          & 0.741          & 0.390          & 0.391          & 0.391          \\
		M2-InternVL3         & 0.000          & 0.000          & 0.000          & 0.298          & \textbf{1.000} & 0.459          \\
		M3-DINOv2            & 0.712          & 0.936          & 0.809          & 0.419          & 0.108          & 0.172          \\
		M3-ResNet-50         & 0.710          & 0.927          & 0.804          & 0.391          & 0.110          & 0.171          \\
		M3-InternVL3         & 0.719          & 0.698          & 0.708          & 0.335          & 0.359          & 0.347          \\
		\bottomrule
	\end{tblr}
\end{table}

The macro values in \cref{tab:binary_routes} are the unweighted means of the two corresponding class-level values reported here.
M1 and the trained M3 routes show high recall for classes A to C (0.927 to 0.960) but low recall for classes D to G (0.098 to 0.110).
M2-DINOv2 provides the most balanced class-level result, with a recall of 0.469 for classes D to G and the highest precision for classes A to C.
M2-InternVL3 assigns every building to classes D to G; its F1 score of 0.459 for these classes therefore reflects the class prevalence (precision = 0.298) rather than discrimination between buildings.
M0 reports the analytical expectation, with precision equal to each class prevalence and recall equal to 0.5.
Source: \texttt{reports/stage3/\allowbreak binary\_prf\_htr\_r2.json}.
